\subsection{Worked Example}

The running example now supports the whole theorem chain with almost no extra notation. The collision block $\{A,B\}$ gives $A_\pi = 2$, so the ambiguity converse forces one tag bit for zero-error identification. The minimal distinguishing family has size $d=3$, so every tag-free zero-error scheme pays at least three queries. The exact block law gives $L_t^\star = t$, and the open-world section uses the same example to show how adding $B$ destroys present collision-freedom in the realized world $\{A,C,D,E\}$.

The following consequences show how the same collision structure governing the ambiguity converse reappears in maintenance locality, architectural centralization, and concrete data-management settings.

\subsection{Maintenance and Localization Consequences}

\paragraph{Scope of this subsection.}
This subsection studies \emph{maintenance or localization} complexity, measured by locations to inspect ($E(\cdot)$), not per-instance identification complexity ($W$) from earlier sections. The metrics are related but distinct: $W$ quantifies online class-identification effort, while $E$ quantifies where constraint logic is distributed in implementations.

\begin{definition}[Error location count]
Let $E(\mathcal{O})$ be the number of locations that must be inspected to find all potential violations of a constraint under observation family $\mathcal{O}$.
\end{definition}
\leanmeta{ \LH{ACS3}, \LH{ACS7}}

\begin{theorem}[Nominal-tag localization]\label{thm:nominal-localization}
$E(\text{nominal-tag}) = O(1)$.
\end{theorem}
\leanmeta{ \LH{ACS7}, \LH{DOM1}, \LH{FORC2}}

\begin{proof}
Under nominal-tag observation, the constraint ``$v$ must be of class $A$'' is satisfied iff $\tau(v) \in \mathrm{subtypes}(A)$. This is determined at a single location: the definition of $\tau(v)$'s class.
\end{proof}

\begin{theorem}[Declared-attribute localization]\label{thm:declared-localization}
$E(\text{attribute-only, declared}) = O(k)$ where $k$ is the number of entity classes.
\end{theorem}
\leanmeta{ \LH{ACS3}}

\begin{proof}
With declared attribute sets, the constraint ``$v$ must satisfy attribute $I$'' requires verifying that each class satisfies all attributes in $I$. For $k$ classes, this yields $O(k)$ locations.
\end{proof}

\begin{theorem}[Attribute-only localization]\label{thm:attribute-localization}
$E(\text{attribute-only}) = \Omega(m)$ where $m$ is the number of query sites.
\end{theorem}
\leanmeta{ \LH{ACS3}, \LH{ACS9}, \LH{COH1}}

\begin{proof}
Under attribute-only observation, each query site independently checks ``does $v$ have attribute $a$?'' with no centralized declaration. For $m$ query sites, each must be inspected. The lower bound is therefore $\Omega(m)$.
\end{proof}

\begin{corollary}[Strict dominance]\label{cor:strict-dominance}
Nominal-tag observation strictly dominates attribute-only: $E(\text{nominal-tag}) = O(1) < \Omega(m) = E(\text{attribute-only})$ for all $m > 1$.
\end{corollary}
\leanmeta{ \LH{ACS4}, \LH{DOM1}, \LH{FORC2}}

\subsection{Information Scattering and Cross-Domain Illustrations}

\begin{definition}[Constraint encoding locations]
Let $I(\mathcal{O}, c)$ be the set of locations where constraint $c$ is encoded under observation family $\mathcal{O}$.
\end{definition}
\leanmeta{ \LH{ACS3}, \LH{ACS6}}

\begin{theorem}[Attribute-only scattering]\label{thm:attribute-scattering}
For attribute-only observation, $|I(\text{attribute-only}, c)| = O(m)$ where $m$ is the number of query sites using constraint $c$.
\end{theorem}
\leanmeta{ \LH{ACS3}}

\begin{proof}
Each attribute query independently encodes the constraint. No shared reference exists. Constraint encodings therefore scale with query sites.
\end{proof}

\begin{theorem}[Nominal-tag centralization]\label{thm:nominal-centralization}
For nominal-tag observation, $|I(\text{nominal-tag}, c)| = O(1)$.
\end{theorem}
\leanmeta{ \LH{ACS6}, \LH{UNIQ1}, \LH{DOM1}}

\begin{proof}
The constraint ``must be of class $A$'' is encoded once in the definition of $A$. All tag checks reference this single definition.
\end{proof}

\begin{corollary}[Maintenance entropy]\label{cor:maintenance-entropy}
Attribute-only observation maximizes maintenance entropy; nominal-tag observation minimizes it.
\end{corollary}
\leanmeta{ \LH{ACS2}, \LH{COH1}, \LH{FORC2}}

\textbf{Databases (columns vs keys).} Rows may collide on non-key columns, so attribute-only identification is insufficient for guaranteed identity. Primary or surrogate keys realize the nominal-tag role and provide constant-cost identity checks \cite{Codd1990}.

\textbf{Model registries and artifact stores.} Structural signatures can collide across distinct artifacts. Registry identifiers and hashes play the same nominal role as tags, while profile inspection alone inherits the ambiguity lower bounds.

\textbf{Biology (phenotype vs genotype).} Distinct species can collide under phenotype, so phenotype-only observation cannot guarantee zero-error identity. DNA barcoding supplies additional naming information that resolves collision blocks in constant-time lookup style \cite{DNABarcoding}.

\textbf{Software runtimes (attribute probes vs runtime identifiers).} Attribute-probe checks recover identity from observed capabilities and pay query cost that scales with inspected structure. Runtime class or type identifiers behave as explicit tags and collapse identity checks to near-constant cost in the model \cite{CPythonDocs,JVMSpec,JavaDocs,TypeScriptDocs,RustDocs}.

These illustrations are not historical causality claims. They show a common resource law: when observable profiles collide, zero-error identification requires additional naming information; otherwise the system pays higher query cost, higher maintenance complexity, and/or nonzero error.

The operational metrics in this section complement the main $L$-, $W$-, and $D$-resource laws. The same collision structure that drives the ambiguity converse also reappears in maintenance burden, error localization, and architectural scattering.
